% Source: graduate-coursework/Prior_Semesters/Fa23/ELCT432/exercises/PA5_Channel_Erasure_Probability.tex
% Description: CDF of Random Variable \(X\)

\documentclass[border=4pt]{standalone}
\usepackage{amsmath}
\usepackage{tikz}

\begin{document}
\begin{tikzpicture}[scale=1.5]
    \draw[->] (0,0) -- (6,0) node[right, font=\small] {\(X\)};
    \draw[->] (0,0) -- (0,6) node[above, font=\small] {Probability};
    \draw[thick] (1,0.625) -- (2,0.625);
    \draw[thick] (2,2.5) -- (3,2.5);
    \draw[thick] (3,4.375) -- (4,4.375);
    \draw[thick] (4,5) -- (5,5);
    
    \foreach \x/\y in {1/0.125, 2/0.5, 3/0.875, 4/1} {
      \fill (\x,\y*5) circle [radius=0.07];
      \node[align=center, below, font=\small] at (\x, \y*5-0.1) {\(\y\)};
    }
    
    \foreach \x/\y in {1/0, 2/1, 3/2, 4/3, 5/4} {
      \node[align=center, below, font=\small] at (\x, -0.1) {\(\y\)};
    }
  \end{tikzpicture}
\end{document}
